% Generated draft: requires mathematical and source review.
\documentclass[11pt]{article}
\usepackage[margin=25mm]{geometry}
\usepackage{fontspec}
\setmainfont{Latin Modern Roman}
\usepackage{amsmath,amssymb,mathtools,unicode-math}
\setmathfont{Latin Modern Math}
\usepackage{enumitem,needspace,xurl,hyperref}
\setlength{\parindent}{0pt}
\setlength{\parskip}{5pt}
\setlength{\emergencystretch}{3em}
\newcommand{\N}{\mathbb{N}}
\newcommand{\Z}{\mathbb{Z}}
\newcommand{\Q}{\mathbb{Q}}
\newcommand{\R}{\mathbb{R}}
\newcommand{\C}{\mathbb{C}}
\newcommand{\F}{\mathbb{F}}
\newcommand{\A}{\mathbb{A}}
\newcommand{\PP}{\mathbb{P}}
\newcommand{\E}{\mathbb{E}}
\newcommand{\eps}{\varepsilon}
\newcommand{\dd}{\,\mathrm{d}}
\begin{document}
\section*{UnsolvedMath ID: 800757}
\textbf{Source problem code:} catalogue entry
\subsection{Definitions and mathematical statement}
Let \(\omega_1\) denote the first uncountable ordinal, and let \([\omega_1]^{<\omega}\) denote the set of its finite subsets. A family \(\mathcal F\subseteq[\omega_1]^{<\omega}\) is hereditary if \(A\in\mathcal F\) and \(B\subseteq A\) imply \(B\in\mathcal F\). For the formulation below, identify each subset \(A\subseteq\omega_1\) with its characteristic function in \(\{0,1\}^{\omega_1}\), equipped with the product topology of the discrete topology on \(\{0,1\}\). Compactness of \(\mathcal F\) means compactness in this topology. This conventional topological interpretation is not explicitly specified in the supplied statement.

The question attributed to D. H. Fremlin is whether there exists a compact hereditary family \(\mathcal F\subseteq[\omega_1]^{<\omega}\) such that
\[
\{\alpha\}\in\mathcal F\qquad(\alpha<\omega_1)
\]
and
\[
\forall A\in[\omega_1]^{<\omega}\ \exists B\subseteq A\quad
\bigl(B\in\mathcal F\ \text{and}\ 2|B|\ge |A|\bigr).
\]
Here the source phrase ``half of its size'' is interpreted as ``at least half'', so the required size is at least \(\lceil |A|/2\rceil\). Heredity then permits choosing \(B\) of exactly that size. The source does not explicitly state the rounding convention for odd cardinalities; this interpretation requires review. For \(A=\varnothing\), take \(B=\varnothing\), which belongs to \(\mathcal F\) by heredity and the singleton condition.
\subsection{Short English statement}
Does there exist a compact hereditary family of finite subsets of \(\omega_1\), containing every singleton, such that every finite set contains a member of the family of at least half its cardinality? This draft uses the characteristic-function product topology and rounds half upward; both conventions require confirmation from the source context.
\subsection{Sources}
\begin{itemize}
\item Supplied local UnsolvedMath registry record: permanent ID \(800757\), exact source problem code \texttt{catalogue entry}, title \emph{Banach Spaces and Ramsey Theory --- Problem 22}. The registry labels the problem \emph{open}; this supplied status has not been independently verified.
\item Dodos and Lopez-Abad, \emph{Banach Spaces and Ramsey Theory: Some Open Problems} (2010), Problem 22. Supplied source URL: \url{https://arxiv.org/abs/1006.2668}. The record reports access on 2026-10-16 and extraction by \texttt{arxiv\_latex}; these are supplied provenance details, not a claim of independent source consultation.
\item The supplied statement attributes the question to D. H. Fremlin, with reference label [Fr2]. No bibliographic entry identifying [Fr2] is supplied, so the original reference cannot be identified precisely from this record.
\end{itemize}
\end{document}
